\documentclass[10pt,a4paper]{amsart}
\usepackage[margin=19mm]{geometry}
\usepackage{amsmath,amssymb,mathtools}
\usepackage[scr=rsfs,cal=euler]{mathalfa}
\usepackage{xfrac}
\usepackage[unicode,hidelinks,bookmarks=false]{hyperref}
\newcommand{\ie}{i.e.\ }
\newcommand{\cf}{cf.\ }
\newcommand{\resp}{respectively\ }
\newcommand{\Subsection}{\subsection}
\providecommand{\nobreakdash}{\nobreak\hbox{-}}
\setlength{\emergencystretch}{2em}
\multlinegap0pt
\usepackage{bm}
\usepackage{bbm}
\usepackage{dsfont}
\usepackage{xspace}
\usepackage{cancel}
\usepackage{mathtools}
\usepackage{amsthm}
\makeatletter
\@ifundefined{theorem}{\newtheorem{theorem}{Theorem}}{}
\makeatother
\title[New recursion formula for the interior polynomial based on n]{New recursion formula for the interior polynomial based on non-expanding sets}
\author{Keiju Kato}
\date{Source: arXiv:2502.19799v3 (CC BY 4.0)}
\begin{document}
\maketitle

\noindent\emph{Source and licence.} New recursion formula for the interior polynomial based on non-expanding sets. Version arXiv:2502.19799v3 (CC BY 4.0), distributed under the Creative Commons Attribution 4.0 International licence, as recorded in the arXiv licence metadata for this exact version and shown on the abstract page https://arxiv.org/abs/2502.19799v3. Full rights notice: licenses/arxiv-2502.19799-CC-BY-4.0.txt.\par\smallskip

\noindent\textit{Editorial scope.} Counted unit: the Theorem whose source label is \texttt{thm:complete} in the article (source0.tex lines 687--692, source label \texttt{thm:complete}), delivered together with the article's own complete original proof of it (source0.tex lines 694--789, one \texttt{proof} environment of 96 lines). No mathematics is written, repaired or re-derived: statement and proof are the article's own text line for line. Required context retained uncounted: combiprec (lines 682--684). Every definition, display and label of those lines is retained unchanged. Omitted from this package and not redistributed: 1--681, 685--686, 693--693, 790--807. Nothing inside a retained excerpt is omitted or altered: every retained body is delivered exactly as the article's lines are written, so no omission receipt of any kind accompanies this package. Citations: the retained excerpts carry no citation command at all, so no bibliography entry of the article is transcribed and no reference is redistributed. Reference adaptation: none was needed and none was made; every reference inside the delivered unit resolves inside it, so no unresolved reference or citation is produced. Printed slips kept literal: no printed word, symbol, hypothesis or number of the counted statement or of its proof was altered. Layout and class: the article's own class is \texttt{amsart} with packages including graphicx, color, float, array, booktabs, amssymb, subcaption, txfonts, amsmath, bm, mathrsfs, overpic, multimedia; the wrapper is the lane's pinned \texttt{amsart} editorial wrapper at 10pt on A4, which restates the theorem-like environments the retained text needs and adds the article's own macro definitions that the retained text needs (none), with extra wrapper packages (bm, bbm, dsfont, xspace, cancel, mathtools). The delivered numbering is the unit's own. Assets: no retained span contains a figure environment, a graphics-inclusion command, a PDF-page inclusion, a table or a TikZ picture; no image file is copied into this package, the package declares no asset dependency and no image file is redistributed with it. Licence: version arXiv:2502.19799v3 of this article is distributed under the Creative Commons Attribution 4.0 International licence, as recorded in the arXiv licence metadata for this exact version and shown on the abstract page https://arxiv.org/abs/2502.19799v3. A full rights notice is delivered at licenses/arxiv-2502.19799-CC-BY-4.0.txt. Characters: every retained excerpt is pure ASCII, so the delivered engine, XeLaTeX with its default Latin Modern text font, reports no missing character.
\begin{equation}\label{combiprec}
I_{K_{m,n}}(x) = \sum_{k=1}^{n} (-1)^{k-1} \binom{n}{k} I_{K_{m,n-k}}(x).
\end{equation}

\begin{theorem}\label{thm:complete}
Let $K_{m,n}$ be the complete bipartite graph. Then the interior polynomial of $K_{m,n}$ is given by
\[
I_{K_{m,n}}(x)=\sum_{j=0}^{\min\{m-1,n-1\}}\binom{m-1}{j}\binom{n-1}{j}x^j.
\]
\end{theorem}

\begin{proof}
First, we may assume $m \le n$ without loss of generality, since the interior polynomial does not depend on the choice of the color classes in the bipartite graph. We use induction on $(m,n)$ ordered lexicographically.

We first establish the base cases.  
For the degenerate case $K_{m,0}$ (i.e., one part is empty), the graph consists of $m$ isolated vertices, so we have
\[
I_{K_{m,0}}(x)=(1-x)^{m-1}.
\]
Using the generalized coefficient convention, where $\binom{-1}{j}=(-1)^j$, this formula can be rewritten as
\[
\sum_{j=0}^{m-1}\binom{m-1}{j}\binom{-1}{j}x^j=\sum_{j=0}^{m-1}\binom{m-1}{j}(-1)^j x^j=(1-x)^{m-1}.
\]
Thus, the formula holds for $K_{m,0}$.

Next, consider the case $m=1$. The graph $K_{1,n}$ is a tree, so we have
\[
I_{K_{1,n}}(x)=\sum_{j=0}^{0}\binom{0}{j}\binom{n-1}{j}x^j=1.
\]

Now, assume that for all pairs $(m',n')$ such that either $m'<m$ or $m'=m$ and $n'<n$, the interior polynomial satisfies
\[
I_{K_{m',n'}}(x)=\sum_{j=0}^{m'-1}\binom{m'-1}{j}\binom{n'-1}{j}x^j.
\]
We proceed to prove the formula for $K_{m,n}$ using the recurrence relation \eqref{combiprec}. Applying the induction hypothesis, we derive
\begin{eqnarray*}
I_{K_{m,n}}(x) &=& \sum_{k=1}^{n} (-1)^{k-1} \binom{n}{k} I_{K_{m,n-k}}(x)\\
&=&\sum_{k=1}^{n}(-1)^{k-1}\binom{n}{k}\left(\sum_{j=0}^{m-1}\binom{m-1}{j}\binom{n-k-1}{j}x^j\right)\\
&=&\sum_{j=0}^{m-1}\binom{m-1}{j}x^j\left(\sum_{k=1}^{n}(-1)^{k-1}\binom{n}{k}\binom{n-k-1}{j}\right).
\end{eqnarray*}

To complete the proof, we need to establish the following identity involving binomial coefficients. For any integer $0 \le j \le n-1$, we have
\begin{equation}\label{eq:binom_identity}
\sum_{k=1}^{n}(-1)^{k-1}\binom{n}{k}\binom{n-k-1}{j}=\binom{n-1}{j}.
\end{equation}

We prove the identity by induction on $n$. First, consider the base case $n=1$. The only possible value of $j$ is $j=0$. In this case, both sides evaluate to 1. Thus, the identity holds for $n=1$.

Now, suppose that for a fixed $n \geq 1$ and for all $0 \le j \le n-1$, the equality
\[
\sum_{k=1}^{n}(-1)^{k-1}\binom{n}{k}\binom{n-k-1}{j} = \binom{n-1}{j}
\]
holds.

For $ n+1$, let $S_{(n+1,j)}$ denote the left-hand side. Then, we compute
\begin{eqnarray*}
S_{(n+1,j)} &=& \sum_{k=1}^{n+1} (-1)^{k-1}\binom{n+1}{k}\binom{n+1-k-1}{j} \\
&=& \sum_{k=1}^{n+1} (-1)^{k-1}\binom{n+1}{k}\binom{n-k}{j}  \\
&=& \sum_{k=1}^{n} (-1)^{k-1}\binom{n+1}{k}\binom{n-k}{j} + (-1)^{n+j}.
\end{eqnarray*}
Using the identity $\binom{n+1}{k} = \binom{n}{k} + \binom{n}{k-1}$, we obtain
\[
S_{(n+1,j)} = \sum_{k=1}^{n} (-1)^{k-1} \left( \binom{n}{k} + \binom{n}{k-1} \right) \binom{n-k}{j} + (-1)^{n+j}.
\]
Splitting the sum, we rewrite it as
\[
S_{(n+1,j)} = S' + S'' + (-1)^{n+j},
\]
where
\[
S' = \sum_{k=1}^{n} (-1)^{k-1} \binom{n}{k} \binom{n-k}{j}, \quad
S'' = \sum_{k=1}^{n} (-1)^{k-1} \binom{n}{k-1} \binom{n-k}{j}.
\]
For $S'$, using the identity $\binom{n-k}{j} = \binom{n-k-1}{j} + \binom{n-k-1}{j-1}$, we obtain
\[
S' = \sum_{k=1}^{n} (-1)^{k-1} \binom{n}{k} \binom{n-k-1}{j}
+ \sum_{k=1}^{n} (-1)^{k-1} \binom{n}{k} \binom{n-k-1}{j-1}.
\]
By the induction hypothesis, we have
\[
S' = \binom{n-1}{j} + \binom{n-1}{j-1} = \binom{n}{j}.
\]
For $S''$, changing variables by setting $l=k-1$, so that $l$ runs from $0$ to $n-1$, we obtain
\begin{eqnarray*}
S'' 
&=& \sum_{l=0}^{n-1} (-1)^{l}\binom{n}{l}\binom{n-l-1}{j} \\
&=&(-1)^{0}\binom{n}{0}\binom{n-0-1}{j}+\sum_{l=1}^{n} (-1)^{l}\binom{n}{l}\binom{n-l-1}{j}-(-1)^{n}\binom{n}{n}\binom{n-n-1}{j} \\
&=&\binom{n-1}{j}-\sum_{l=1}^{n} (-1)^{l-1}\binom{n}{l}\binom{n-l-1}{j}-(-1)^n\binom{-1}{j}
\end{eqnarray*}
By the induction hypothesis, we have
\[
S''=\binom{n-1}{j}-\binom{n-1}{j}-(-1)^{n+j}=-(-1)^{n+j}
\]
Combining all terms, we obtain
\[
S_{(n+1, j)} =S' + S'' +(-1)^{n+j} 
= \binom{n}{j}-(-1)^{n+j} +(-1)^{n+j} 
= \binom{n}{j}.
\]
Thus, we have established the identity \eqref{eq:binom_identity}, completing the proof.

From \eqref{eq:binom_identity}, we obtain 
\[
I_{K_{m,n}}(x)=\sum_{j=0}^{m-1}\binom{m-1}{j}\binom{n-1}{j}x^j.
\]
Using induction on $(m,n)$ ordered lexicographically, we conclude that the stated formula for the interior polynomial of $K_{m,n}$ holds for all $ m \le n$.
\end{proof}

\end{document}
